\chapter{Resultados complementarios da avaliación de Tesseract}
\label{chapapd:resultados-tesseract}

Neste apéndice recóllense exemplos adicionais dos resultados obtidos durante
a avaliación do modelo estándar de Tesseract. O obxectivo é complementar os
casos mostrados no corpo principal mediante fragmentos máis amplos das saídas
xeradas e a súa comparación directa coas follas de tratamento orixinais.

A avaliación realizouse sobre as \textbf{11 follas escaneadas}, empregadas como
escenario de referencia ao representar as condicións de entrada máis
favorables. Dado que xa neste conxunto se observaron erros importantes no
recoñecemento e na conservación da estrutura da táboa de doses, non se
considerou necesario ampliar a proba do modelo estándar a fotografías tomadas
co dispositivo móbil.

Os exemplos seguintes non pretenden recoller todos os erros observados, senón
mostrar distintos comportamentos representativos: a degradación dunha pauta
extensa, a omisión da táboa de doses, os erros en valores numéricos e,
finalmente, un caso no que o recoñecemento da pauta foi considerablemente
mellor.


\section{Degradación da táboa de doses}
\label{secapd:tesseract-f01}

A folla F01 contén unha pauta extensa distribuída ao longo de varias semanas.
Neste caso, Tesseract recupera correctamente unha parte das primeiras entradas,
pero a saída vaise degradando e aparecen modificacións nas datas e nas
fraccións, ademais de caracteres que non forman parte da información da
táboa.

A Figura~\ref{fig:apendice-tesseract-f01} mostra o fragmento correspondente da
folla orixinal, mentres que o
Listing~\ref{lst:apendice-tesseract-f01} recolle a saída xerada por Tesseract.

\begin{figure}[htbp]
    \centering
    % Recortar a táboa da folla F01.
    % Interesa que se vexan especialmente as últimas filas,
    % nas que a saída de Tesseract presenta máis ruído.
    \includegraphics[width=0.93\textwidth]
    {imaxes/apendices/tesseract-f01-orixinal.jpg}
    \caption{Fragmento da táboa de doses da folla F01 empregada na
    avaliación do modelo estándar de Tesseract.}
    \label{fig:apendice-tesseract-f01}
\end{figure}

\begin{lstlisting}[
    frame=single,
    rulecolor=\color{gray!60},
    backgroundcolor=\color{gray!5},
    basicstyle=\ttfamily\scriptsize,
    caption={Fragmento da saída de Tesseract correspondente á
    Figura~\ref{fig:apendice-tesseract-f01}},
    label={lst:apendice-tesseract-f01},
    captionpos=t,
    columns=fullflexible,
    keepspaces=true
]
LUNES MARTES MIERCOLES JUEVES VIERNES SABADO DOMINGO
08 SEP. 1/2 09 SEP 1/4 |l0SEP 1/2 |11 SEP 1/2 |12SEP 1/2 13SEP 1/4 |14SEP 1/2
15 SEP. 1/2 lI6SEP 1/4 I17SEP 1/2 |I8 SEP 1/2 |19SEP 1/2 20SEP 1/4 121 SEP 172
22 SEP 1/2 (23 SEP 114 |2SEP 12 |125SEP 1/2 26SEP 1/2 |27 SEP 1/4 28 SEP 1/2
29 SEP" 1/2 |30 SEP 1/4 |01 12 [02 12 10 12 04 1/4 [05 1/2
al O a a d a d
06 12 [07 1/14 08 12 [09 1/2 |10 12 | 1 14 |2 1/2
OCT OCT OCT OCT OCT OCT OCT
\end{lstlisting}

A comparación permite observar que unha parte do contido segue sendo
recoñecible, pero aparecen modificacións como a perda dos separadores das
fraccións, caracteres adicionais e alteracións nos días. Tamén se introduce
unha liña de ruído sen correspondencia coa información da pauta. Como
resultado, a secuencia completa deixa de poder interpretarse de maneira
fiable mediante un procesamento textual directo.


\section{Omisión da pauta diaria}
\label{secapd:tesseract-f02}

Un comportamento diferente apareceu na folla F02. Neste caso, Tesseract
consegue recoñecer diferentes elementos do documento, incluído o encabezado
da sección de doses, pero non recupera o contido da táboa situada baixo el.

A Figura~\ref{fig:apendice-tesseract-f02} mostra a zona correspondente da
folla orixinal. O Listing~\ref{lst:apendice-tesseract-f02} permite comprobar
que, tras recoñecer o título da táboa, a saída pasa directamente á información
da próxima visita.

\begin{figure}[htbp]
    \centering
    % Recortar desde "DOSIS DIARIA DE ANTICOAGULANTE"
    % ata xusto antes de "Próxima visita".
    % Debe verse claramente a pauta presente no documento orixinal.
    \includegraphics[width=0.93\textwidth]
    {imaxes/apendices/tesseract-f02-orixinal.jpg}
    \caption{Táboa de doses da folla F02 que non foi recuperada na saída
    textual de Tesseract.}
    \label{fig:apendice-tesseract-f02}
\end{figure}

\begin{lstlisting}[
    frame=single,
    rulecolor=\color{gray!60},
    backgroundcolor=\color{gray!5},
    basicstyle=\ttfamily\scriptsize,
    caption={Fragmento da saída de Tesseract correspondente á
    Figura~\ref{fig:apendice-tesseract-f02}},
    label={lst:apendice-tesseract-f02},
    captionpos=t,
    columns=fullflexible,
    keepspaces=true
]
Fármaco oral.:Sintrom 1 mg
Dosis semanal: 4 mg (1,0,1,0,1,0,1) Dosis inyecta:

DOSIS DIARIA DE ANTICOAGULANTE

Próxima visita: martes, 29/04/2025 Turno: Mañanas - Num: 5
Centro visita:
Comentarios.
C.S. TABOADA
\end{lstlisting}

Neste caso, o problema non se limita á reorganización ou á modificación dos
caracteres: a información correspondente ás datas e doses non aparece na
saída. Polo tanto, non sería posible recuperar a pauta unicamente mediante a
aplicación posterior de regras de reorganización ou expresións regulares.


\section{Modificación de valores numéricos}
\label{secapd:tesseract-f03}

Os erros non se limitaron á táboa de doses. Na folla F03 tamén se observaron
modificacións en campos numéricos do resumo das últimas visitas. A
Figura~\ref{fig:apendice-tesseract-f03} mostra o fragmento orixinal empregado
para a comparación e o Listing~\ref{lst:apendice-tesseract-f03} presenta a
saída obtida.

\begin{figure}[htbp]
    \centering
    % Recortar "RESUMEN DE LAS ULTIMAS VISITAS"
    % incluíndo as filas que aparecen no listing.
    \includegraphics[width=0.93\textwidth]
    {imaxes/apendices/tesseract-f03-orixinal.jpg}
    \caption{Fragmento do resumo das últimas visitas da folla F03.}
    \label{fig:apendice-tesseract-f03}
\end{figure}

\begin{lstlisting}[
    frame=single,
    rulecolor=\color{gray!60},
    backgroundcolor=\color{gray!5},
    basicstyle=\ttfamily\scriptsize,
    caption={Fragmento da saída de Tesseract correspondente á
    Figura~\ref{fig:apendice-tesseract-f03}},
    label={lst:apendice-tesseract-f03},
    captionpos=t,
    columns=fullflexible,
    keepspaces=true
]
15/01/2025 8  Sintrom4mg 35 me 16/01/2025
07/01/2025 29 Sintrom4mg 4 mg 16/01/2025
30/12/2024 54 Sintrom4mg 4 mg 07/01/2025
23/12/2024 44 Sintrom4mg 53 mg 30/12/2024
16/12/2024 3,1 Sintrom4mg 5,6 mg 24/12/2024
\end{lstlisting}

A comparación coa imaxe orixinal permite comprobar que varios valores
numéricos son modificados durante o recoñecemento. Este tipo de erro resulta
especialmente relevante porque a saída pode continuar tendo a aparencia dun
valor numérico válido, aínda que xa non coincida co presente no documento.
Por tanto, a detección posterior do erro non pode basearse unicamente en
comprobar que o resultado teña un formato numérico.


\section{Caso cun mellor recoñecemento da pauta}
\label{secapd:tesseract-f05}

Os resultados non presentaron o mesmo nivel de degradación en todas as
follas. Na F05, que contén unha pauta máis curta, Tesseract conseguiu conservar
con bastante fidelidade a secuencia principal de datas e doses.

Inclúese este exemplo para mostrar a variabilidade observada entre os
documentos e evitar limitar a análise unicamente aos casos nos que se
produciron os erros máis evidentes. A
Figura~\ref{fig:apendice-tesseract-f05} mostra a pauta orixinal e o
Listing~\ref{lst:apendice-tesseract-f05} a saída correspondente.

\begin{figure}[htbp]
    \centering
    % Recortar a fila completa de LUNES a DOMINGO,
    % do 16 DIC ao 22 DIC, incluíndo as doses e "NO TOMAR".
    \includegraphics[width=0.93\textwidth]
    {imaxes/apendices/tesseract-f05-orixinal.jpg}
    \caption{Táboa de doses da folla F05, correspondente a un dos casos
    cun mellor recoñecemento por parte do modelo estándar de Tesseract.}
    \label{fig:apendice-tesseract-f05}
\end{figure}

\begin{lstlisting}[
    frame=single,
    rulecolor=\color{gray!60},
    backgroundcolor=\color{gray!5},
    basicstyle=\ttfamily\scriptsize,
    caption={Saída de Tesseract correspondente á
    Figura~\ref{fig:apendice-tesseract-f05}},
    label={lst:apendice-tesseract-f05},
    captionpos=t,
    columns=fullflexible,
    keepspaces=true
]
| LUNES MARTES MIERCOLES JUEVES VIERNES SABADO DOMINGO
16 DIC 0 17 DIC 1/4 18 DIC 1/4 19 DIC 1/4 20 DIC 1/4 21 DIC 0 22 DIC 1/4
No Tomar 8 a a a No Tomar [(
\end{lstlisting}

Neste caso, a asociación principal entre os días e as doses mantense de forma
moito máis clara ca nos exemplos anteriores. A saída segue introducindo
caracteres adicionais na liña inferior, pero a pauta principal resulta
recoñecible.

Este resultado mostra que a limitación do modelo estándar non consiste en que
sexa incapaz de recuperar calquera táboa deste tipo, senón na falta de
consistencia entre documentos. Unha solución destinada a reconstruír
automaticamente a pauta de tratamento non pode asumir que unha estrutura que
se conserva nunha folla vaia manterse nas restantes.


\section{Síntese dos exemplos}
\label{secapd:tesseract-sintese}

Os exemplos anteriores mostran que o comportamento do modelo estándar varía
considerablemente entre as distintas follas escaneadas. Nalgúns documentos
consérvase unha parte importante da pauta, mentres que noutros aparecen
modificacións nas fraccións e nas datas, ruído adicional ou mesmo a omisión
completa do contido da táboa.

Esta variabilidade é especialmente problemática neste proxecto, xa que a
extracción require manter de maneira fiable a relación entre cada data e a súa
dose. A presenza destas limitacións nas propias follas escaneadas foi
suficiente para non continuar a avaliación do modelo estándar en condicións de
captura menos favorables e motivou a experimentación posterior cun modelo de
Tesseract adestrado especificamente para este tipo de documentos.